\chapter{高等数学：数列极限}

\section{数列极限}

\begin{problemblock}
\textbf{例1.} “对任意的 \(\varepsilon>0\)，只有有限个
\[
x_n\notin[a-\varepsilon,a+\varepsilon]
\]
”是“\(\displaystyle\lim_{n\to\infty}x_n=a\)”的（\quad）

A. 充分不必要条件 \qquad B. 必要不充分条件

C. 充要条件 \qquad D. 既不充分也不必要条件
\end{problemblock}

\begin{solutionblock}
\analysis{本题考查数列极限定义的等价表述。定义说：对任意 \(\varepsilon>0\)，存在 \(N\)，使 \(n>N\) 时 \(|x_n-a|<\varepsilon\)。这等价于“落在邻域外的项至多有限个”。}

若 \(\lim x_n=a\)，则对任意 \(\varepsilon>0\)，存在 \(N\)，当 \(n>N\) 时
\[
x_n\in[a-\varepsilon,a+\varepsilon].
\]
所以不在该区间内的项只能出现在 \(x_1,\dots,x_N\) 中，至多有限个。

反过来，若对任意 \(\varepsilon>0\)，只有有限个 \(x_n\notin[a-\varepsilon,a+\varepsilon]\)，则这些例外项有最大下标 \(N\)。于是当 \(n>N\) 时，
\[
x_n\in[a-\varepsilon,a+\varepsilon],
\]
即 \(|x_n-a|\le\varepsilon\)。把 \(\varepsilon\) 换成 \(\varepsilon/2\) 即得标准定义中的严格小于。

\answer{C}
\examnote{“有限个例外项”与“充分靠后的所有项”是数列极限定义中常用的等价语言。}
\end{solutionblock}

\begin{problemblock}
\textbf{例2.} 关于正项数列 \(\{a_n\}\)，下列说法正确的有几项？
\begin{enumerate}[label=\arabic*., leftmargin=2.2em]
\item 若 \(\left\{\dfrac{a_n}{\ln a_n}\right\}\) 收敛，则 \(\{a_n\}\) 一定收敛；
\item 若 \(\{a_n\}\) 收敛，则 \(\left\{\dfrac{a_n}{\ln a_n}\right\}\) 一定收敛；
\item 若 \(\{\arctan a_n\}\) 发散，则 \(\{a_n\}\) 一定不单调；
\item 若 \(\{a_n\}\) 发散，则 \(\left\{\dfrac{2a_n}{a_n^2+1}\right\}\) 一定发散；
\item 若 \(\{a_n\}\) 单调增加且 \(\{a_n\ln a_n\}\) 收敛到 \(0\)，则 \(\{a_n\}\) 必收敛；
\item 若 \(\{a_n\}\) 单调增加且 \(\left\{(1+a_n)^{1/a_n}\right\}\) 收敛到 \(1\)，则 \(\{a_n\}\) 必收敛。
\end{enumerate}
\end{problemblock}

\begin{solutionblock}
\analysis{本题是“连续函数作用于数列”与“反推数列收敛”的综合判断。考研中这类题要多用反例。}

①错误。函数
\[
\phi(t)=\frac{t}{\ln t}
\]
在 \((1,+\infty)\) 上不是一一对应：在 \((1,e)\) 上递减，在 \((e,+\infty)\) 上递增。可取两个不同正数 \(u,v>1\)，使 \(\phi(u)=\phi(v)\)，令 \(a_n\) 在 \(u,v\) 间交替，则 \(\phi(a_n)\) 收敛而 \(a_n\) 不收敛。

②错误。若 \(a_n\to1\)，则 \(\ln a_n\to0\)，\(\frac{a_n}{\ln a_n}\) 不一定收敛。例如 \(a_n=1+\frac1n\)，则
\[
\frac{a_n}{\ln a_n}\to+\infty.
\]

③正确。若 \(\{a_n\}\) 单调且 \(a_n>0\)，则 \(\{a_n\}\) 要么收敛，要么趋于 \(+\infty\)。由于 \(\arctan x\) 单调有界，\(\{\arctan a_n\}\) 必收敛。因此若 \(\{\arctan a_n\}\) 发散，\(\{a_n\}\) 一定不单调。

④错误。取 \(a_n=n\)，则 \(\{a_n\}\) 发散，但
\[
\frac{2a_n}{a_n^2+1}=\frac{2n}{n^2+1}\to0.
\]

⑤正确。若 \(\{a_n\}\) 单调增加，则其极限存在于扩展实数中。若 \(a_n\to+\infty\)，则 \(a_n\ln a_n\to+\infty\)，矛盾；若 \(a_n\to L>0\)，由连续性可得 \(L\ln L=0\)，从而 \(L=1\)。总之 \(\{a_n\}\) 收敛。

⑥错误。取 \(a_n=n\)，则 \(\{a_n\}\) 单调增加且发散，而
\[
(1+n)^{1/n}\to1.
\]

正确的是③⑤，共 2 项。

\answer{2 项}
\examnote{“函数值收敛”一般不能反推“自变量数列收敛”，除非函数在相关范围内有单调一一性。}
\end{solutionblock}

\begin{problemblock}
\textbf{例3.} 已知数列 \(\{x_n\}\)，以下结论：
\[
\begin{array}{ll}
\text{① 当 }\lim\limits_{n\to\infty}\cos(\arctan x_n)\text{ 存在时，}\lim\limits_{n\to\infty}x_n\text{ 存在， }-\infty<x_n<+\infty;\\
\text{② 当 }\lim\limits_{n\to\infty}\cos(\tan x_n)\text{ 存在时，}\lim\limits_{n\to\infty}x_n\text{ 存在， }-\frac\pi2<x_n<\frac\pi2;\\
\text{③ 当 }\lim\limits_{n\to\infty}\sin(\arctan x_n)\text{ 存在时，}\lim\limits_{n\to\infty}x_n\text{ 存在， }-\infty<x_n<+\infty;\\
\text{④ 当 }\lim\limits_{n\to\infty}\sin(\tan x_n)\text{ 存在时，}\lim\limits_{n\to\infty}x_n\text{ 存在， }-\frac\pi2<x_n<\frac\pi2.
\end{array}
\]
正确结论的个数为（\quad）

A. 0 \qquad B. 1 \qquad C. 2 \qquad D. 3
\end{problemblock}

\begin{solutionblock}
\analysis{这类题本质是问函数复合后收敛能否推出原数列收敛。关键看外层函数是否一一，是否会丢失符号或周期信息。}

①错误。因为
\[
\cos(\arctan x)=\frac1{\sqrt{1+x^2}},
\]
只反映 \(|x|\)，不反映符号。取
\[
x_n=(-1)^n,
\]
则 \(\cos(\arctan x_n)=1/\sqrt2\) 收敛，但 \(x_n\) 不收敛。

②错误。取 \(x_n=(-1)^n c\)，其中 \(0<c<\frac\pi2\)。则 \(\tan x_n\) 在 \(\pm\tan c\) 间交替，而
\[
\cos(\tan x_n)=\cos(\tan c)
\]
收敛，但 \(x_n\) 不收敛。

③错误。虽然
\[
\sin(\arctan x)=\frac{x}{\sqrt{1+x^2}}
\]
在有限范围内单调，但它的值可以趋于 \(1\) 而 \(x_n\to+\infty\)。取 \(x_n=n\)，则 \(\sin(\arctan x_n)\to1\)，但 \(x_n\) 不收敛于有限实数。

④错误。由于 \(\sin u\) 不是一一函数，可取不同的 \(u,v\) 使 \(\sin u=\sin v\)，再令 \(x_n=\arctan u,\arctan v\) 交替，即可使 \(\sin(\tan x_n)\) 收敛而 \(x_n\) 不收敛。

\answer{A}
\examnote{复合函数收敛反推原数列收敛，要警惕两类失真：偶函数丢符号，周期函数丢位置。}
\end{solutionblock}

\begin{problemblock}
\textbf{例4.} 设数列 \(\{a_n\}\) 满足 \(\displaystyle\lim_{n\to\infty}a_n=0\)。

(1) 若 \(a_1>0,a_2<0\)，证明 \(\{a_n\}\) 有最大值，有最小值。

(2) 若 \(a_1>0,a_2=0\)，证明 \(\{a_n\}\) 有最大值，有最小值。
\end{problemblock}

\begin{solutionblock}
\analysis{收敛到 0 的数列最终会落入任意小邻域。若前面已经有正项或负项，就能把最大值、最小值限制在有限项中取得。}

(1) 因为 \(a_n\to0\)，取
\[
\varepsilon=\min\{a_1,-a_2\}>0.
\]
存在 \(N\)，使 \(n>N\) 时
\[
|a_n|<\varepsilon.
\]
于是当 \(n>N\) 时，
\[
a_2<-\varepsilon<a_n<\varepsilon<a_1.
\]
所以最大值和最小值都只能在有限集合
\[
\{a_1,a_2,\dots,a_N\}
\]
中取得。有限个实数必有最大值和最小值，故结论成立。

(2) 因为 \(a_1>0,a_2=0\)，取 \(\varepsilon=a_1\)。存在 \(N\)，使 \(n>N\) 时
\[
|a_n|<a_1.
\]
所以最大值在有限集合 \(\{a_1,\dots,a_N\}\) 中取得。

对最小值，若所有 \(a_n\ge0\)，则 \(a_2=0\) 就是最小值；若存在某项 \(a_k<0\)，再取一个负项 \(a_k\)，用与(1)相同的方法可知所有足够靠后的项都大于 \(a_k\)，最小值只能在有限项中取得。故最小值存在。

\answer{(1)、(2) 均成立：数列都有最大值，也都有最小值。}
\examnote{“收敛数列有界”不等于“一定有最大最小值”；本题要利用给定的正项、负项或零项把极值锁在有限项中。}
\end{solutionblock}

\begin{problemblock}
\textbf{例5.} 已知
\[
a_n=(-1)^n\sqrt2+\sqrt n,\qquad n=1,2,\dots
\]
则 \(a_n\)（\quad）

A. 既有最大值，也有最小值 \qquad B. 有最大值，没有最小值

C. 没有最大值，有最小值 \qquad D. 既没有最大值，也没有最小值
\end{problemblock}

\begin{solutionblock}
\analysis{本题看似有振荡项，但主项 \(\sqrt n\to+\infty\)，所以整体无上界。最小值需要检查有限前项。}

因为
\[
a_n=\sqrt n+(-1)^n\sqrt2\to+\infty
\]
沿总体趋势无上界，所以没有最大值。

当 \(n\) 为偶数时，
\[
a_n=\sqrt n+\sqrt2>0.
\]
当 \(n\) 为奇数时，
\[
a_n=\sqrt n-\sqrt2.
\]
奇数项中 \(n=1\) 时
\[
a_1=1-\sqrt2<0,
\]
而 \(n=3,5,\dots\) 时 \(\sqrt n-\sqrt2>0\)。因此最小值为
\[
a_1=1-\sqrt2.
\]

\answer{C}
\examnote{“趋于 \(+\infty\)”保证没有最大值，但最小值仍可能由前面有限项取得。}
\end{solutionblock}

\begin{problemblock}
\textbf{例6.} 设数列 \(\{x_n\}\) 满足
\[
x_1=1,\qquad x_n=x_{n+1}+2\ln(1+x_{n+1})\quad(n=1,2,\dots).
\]
(1) 证明 \(\displaystyle\lim_{n\to\infty}x_n\) 存在，并求其值；
(2) 求 \(\displaystyle\lim_{n\to\infty}2^n x_n\)。
\end{problemblock}

\begin{solutionblock}
\analysis{递推式由 \(x_{n+1}\) 反过来定义 \(x_n\)。先证明 \(x_n>0\) 且递减，再由极限方程求极限。第二问看相邻项比值。}

设
\[
\phi(t)=t+2\ln(1+t)\quad(t>-1).
\]
则
\[
\phi'(t)=1+\frac2{1+t}>0,
\]
所以 \(\phi\) 严格递增。由 \(x_1=1>0\) 及 \(\phi(0)=0\)，可递推得 \(x_{n+1}>0\)。又
\[
x_n=x_{n+1}+2\ln(1+x_{n+1})>x_{n+1},
\]
故 \(\{x_n\}\) 单调递减且有下界 0，因此收敛。

设极限为 \(L\ge0\)，令 \(n\to\infty\) 得
\[
L=L+2\ln(1+L),
\]
故
\[
\ln(1+L)=0,\qquad L=0.
\]

第二问，由
\[
x_n=x_{n+1}+2\ln(1+x_{n+1})
\]
且 \(x_{n+1}\to0\)，得
\[
x_n=3x_{n+1}+O(x_{n+1}^2),
\]
所以
\[
\frac{x_{n+1}}{x_n}\to\frac13.
\]
取 \(q\) 满足 \(\frac13<q<\frac12\)，则充分大时 \(x_{n+1}<q x_n\)，从而
\[
x_n\le Cq^n.
\]
于是
\[
0\le2^n x_n\le C(2q)^n\to0.
\]

\answer{(1) 极限为 \(0\)；(2) \(\displaystyle 0\)。}
\examnote{递推极限题常用“先证收敛，再代入极限方程”；带 \(2^n\) 的第二问则看相邻项比值。}
\end{solutionblock}

\begin{problemblock}
\textbf{例7.} 设数列 \(\{x_n\}\) 满足
\[
x_1>0,\qquad x_{n+1}=x_n-\ln(1+x_n).
\]
(1) 证明 \(\{x_n\}\) 收敛，并求其极限；

(2) 设 \(x_1=y_1=\frac12,\ y_{n+1}=y_n^2\)，证明
\[
\lim_{n\to\infty}\frac{x_n}{y_n}=0.
\]
\end{problemblock}

\begin{solutionblock}
\analysis{第一问用 \(\ln(1+x)<x\) 保正，用递减有界证收敛。第二问要比较两个二次级递推的系数。}

当 \(x>0\) 时，
\[
0<\ln(1+x)<x,
\]
所以
\[
0<x_{n+1}=x_n-\ln(1+x_n)<x_n.
\]
故 \(\{x_n\}\) 单调递减且有下界 0，极限存在。设极限为 \(L\ge0\)，则
\[
L=L-\ln(1+L),
\]
从而 \(L=0\)。

第二问中，利用
\[
\ln(1+x)>x-\frac{x^2}{2}\qquad(x>0)
\]
可得
\[
x_{n+1}=x_n-\ln(1+x_n)<\frac{x_n^2}{2}.
\]
又 \(y_{n+1}=y_n^2\)，令
\[
r_n=\frac{x_n}{y_n}.
\]
则
\[
r_{n+1}=\frac{x_{n+1}}{y_{n+1}}<\frac12\left(\frac{x_n}{y_n}\right)^2=\frac12 r_n^2.
\]
由于 \(r_1=1\)，所以 \(r_2<1/2\)，并递推得到 \(0<r_n\to0\)。因此
\[
\lim_{n\to\infty}\frac{x_n}{y_n}=0.
\]

\answer{(1) \(\displaystyle\lim_{n\to\infty}x_n=0\)；(2) \(\displaystyle\lim_{n\to\infty}\frac{x_n}{y_n}=0\)。}
\method{方法二（渐近比较）}
第二问还可以取倒数或取对数比较。由 \(x_{n+1}<x_n^2/2\) 递推得 \(\ln x_{n+1}<2\ln x_n-\ln2\)，而 \(y_n\) 满足完全平方递推；从对数层面可看出 \(x_n\) 比 \(y_n\) 多出连续的 \(1/2\) 系数损失，因此 \(x_n/y_n\to0\)。
\examnote{本题第二问不能只说二者都趋于 0；要比较衰减速度，关键不等式是 \(x-\ln(1+x)<x^2/2\)。}
\end{solutionblock}

\begin{problemblock}
\textbf{例8.} 设数列 \(\{x_n\},\{y_n\}\) 满足
\[
x_1=y_1=\frac12,\qquad x_{n+1}=\sin x_n,\qquad y_{n+1}=y_n^2.
\]
则当 \(n\to\infty\) 时（\quad）

A. \(x_n\) 是 \(y_n\) 的高阶无穷小 \qquad
B. \(y_n\) 是 \(x_n\) 的高阶无穷小

C. \(x_n\) 是 \(y_n\) 的等价无穷小 \qquad
D. \(x_n\) 是 \(y_n\) 的同阶但非等价无穷小
\end{problemblock}

\begin{solutionblock}
\analysis{\(y_n\) 每次平方，衰减极快；\(x_{n+1}=\sin x_n\) 只比 \(x_n\) 略小，衰减很慢。}

由递推可得
\[
y_n=\left(\frac12\right)^{2^{n-1}}.
\]
另一方面，因 \(0<x_n\le1/2\)，且在该范围内 \(\sin x>\frac{x}{2}\)，所以
\[
x_{n+1}>\frac{x_n}{2}.
\]
递推得
\[
x_n>\left(\frac12\right)^n.
\]
于是
\[
0<\frac{y_n}{x_n}
<\frac{2^{-2^{n-1}}}{2^{-n}}
=2^{n-2^{n-1}}\to0.
\]
所以
\[
y_n=o(x_n),
\]
即 \(y_n\) 是 \(x_n\) 的高阶无穷小。

\answer{B}
\examnote{遇到递推比较阶数，先找显式项或粗下界。平方递推通常是超指数衰减。}
\end{solutionblock}

\begin{problemblock}
\textbf{例9.} 设 \(\{x_n\},\{y_n\},\{z_n\}\) 为数列，则以下说法：
\[
\begin{array}{l}
\text{① 若 }x_n\le y_n\le z_n,\text{ 且 }\lim x_n=\lim z_n,\text{ 则 }\{y_n\}\text{ 收敛；}\\
\text{② 若 }x_n\le y_n\le z_n,\text{ 且 }\lim (z_n-x_n)=0,\text{ 则 }\{y_n\}\text{ 收敛；}\\
\text{③ 若 }x_n\le a\le z_n,\text{ 且 }\lim (z_n-x_n)=0,\text{ 则 }\{x_n\}\text{ 和 }\{z_n\}\text{ 收敛；}\\
\text{④ 若 }x_n\le y_n\le z_n,\text{ 且 }\{y_n\}\text{ 收敛，}\lim (z_n-x_n)=0,\text{ 则 }\{x_n\}\text{ 和 }\{z_n\}\text{ 收敛。}
\end{array}
\]
正确的个数是（\quad）

A. 1 个 \qquad B. 2 个 \qquad C. 3 个 \qquad D. 4 个
\end{problemblock}

\begin{solutionblock}
\analysis{夹逼定理要求两端有同一极限；仅有 \(z_n-x_n\to0\) 不足以保证中间数列收敛，因为两端可能一起振荡。}

①正确。这是夹逼定理。

②错误。取
\[
x_n=(-1)^n,\qquad z_n=(-1)^n+\frac1n,\qquad y_n=(-1)^n+\frac1{2n}.
\]
则 \(x_n\le y_n\le z_n\)，且 \(z_n-x_n=1/n\to0\)，但 \(y_n\) 不收敛。

③正确。由
\[
x_n\le a\le z_n
\]
得
\[
0\le a-x_n\le z_n-x_n,\qquad 0\le z_n-a\le z_n-x_n.
\]
因为 \(z_n-x_n\to0\)，所以
\[
x_n\to a,\qquad z_n\to a.
\]

④正确。设 \(y_n\to L\)。由
\[
0\le y_n-x_n\le z_n-x_n,\qquad 0\le z_n-y_n\le z_n-x_n
\]
可得
\[
x_n-y_n\to0,\qquad z_n-y_n\to0.
\]
所以
\[
x_n\to L,\qquad z_n\to L.
\]

正确的是①③④，共 3 个。

\answer{C}
\examnote{“宽度趋于 0”只能说明三个数列彼此接近，不能说明它们有共同固定目标；还需要一个锚点。}
\end{solutionblock}

\begin{problemblock}
\textbf{例10.} 计算定积分
\[
\int_0^\pi\frac{1}{1+\cos^2x}\,dx,\qquad
\int_0^\pi\frac{\sin^2x}{1+\cos^2x}\,dx,
\]
并求极限
\[
\lim_{x\to+\infty}
\frac{\displaystyle\int_0^x\frac{\sin^2t}{1+\cos^2t}\,dt}
{\displaystyle\int_0^x\frac{1}{1+\cos^2t}\,dt}.
\]
\end{problemblock}

\begin{solutionblock}
\analysis{两个被积函数都是以 \(\pi\) 为周期的函数。先算一个周期上的积分，再用周期函数平均值求极限。}

记
\[
I_1=\int_0^\pi\frac{dx}{1+\cos^2x}.
\]
利用对称性，
\[
I_1=2\int_0^{\pi/2}\frac{dx}{1+\cos^2x}.
\]
令 \(t=\tan x\)，则
\[
\cos^2x=\frac1{1+t^2},\qquad dx=\frac{dt}{1+t^2}.
\]
所以
\[
\int_0^{\pi/2}\frac{dx}{1+\cos^2x}
=\int_0^{+\infty}\frac{dt}{t^2+2}
=\frac{\pi}{2\sqrt2}.
\]
故
\[
I_1=\frac{\pi}{\sqrt2}.
\]

再记
\[
I_2=\int_0^\pi\frac{\sin^2x}{1+\cos^2x}\,dx.
\]
因为
\[
\frac{\sin^2x}{1+\cos^2x}
=\frac{1-\cos^2x}{1+\cos^2x}
=-1+\frac{2}{1+\cos^2x},
\]
所以
\[
I_2=-\pi+2I_1
=-\pi+\frac{2\pi}{\sqrt2}
=\pi(\sqrt2-1).
\]

若 \(p(t),q(t)\) 为正的 \(\pi\)-周期函数，则
\[
\frac{\int_0^x p(t)\,dt}{\int_0^x q(t)\,dt}
\to
\frac{\int_0^\pi p(t)\,dt}{\int_0^\pi q(t)\,dt}.
\]
因此所求极限为
\[
\frac{I_2}{I_1}
=\frac{\pi(\sqrt2-1)}{\pi/\sqrt2}
=2-\sqrt2.
\]

\answer{\(\displaystyle I_1=\frac{\pi}{\sqrt2},\quad I_2=\pi(\sqrt2-1),\quad\) 极限为 \(2-\sqrt2\)。}
\examnote{周期积分比值的极限就是“一个周期平均值之比”，这是处理此类极限的快捷方法。}
\end{solutionblock}

\begin{problemblock}
\textbf{例11.} (1) 证明
\[
\frac1{2n}<\frac12\cdot\frac34\cdots\frac{2n-1}{2n}<\frac1{\sqrt{2n+1}}.
\]
(2) 求极限
\[
\lim_{n\to\infty}\sqrt[n]{\frac12\cdot\frac34\cdots\frac{2n-1}{2n}}.
\]
\end{problemblock}

\begin{solutionblock}
\analysis{先用不等式夹住 Wallis 型乘积，再对第 \(n\) 次根取极限。}

记
\[
P_n=\prod_{k=1}^n\frac{2k-1}{2k}.
\]
由 Wallis 不等式可得
\[
\frac1{2n}<P_n<\frac1{\sqrt{2n+1}}.
\]
也可用归纳法证明：左边由 \(P_n>\frac1{2n}\) 递推验证；右边等价于
\[
P_n^2<\frac1{2n+1},
\]
利用
\[
P_n^2=\prod_{k=1}^n\left(\frac{2k-1}{2k}\right)^2
<\prod_{k=1}^n\frac{2k-1}{2k+1}
=\frac1{2n+1}.
\]

对(2)，由(1)得
\[
\sqrt[n]{\frac1{2n}}<\sqrt[n]{P_n}<\sqrt[n]{\frac1{\sqrt{2n+1}}}.
\]
两端均趋于 1，因此夹逼得
\[
\lim_{n\to\infty}\sqrt[n]{P_n}=1.
\]

\answer{极限为 \(1\)。}
\examnote{遇到乘积的 \(n\) 次根，常用“粗阶估计”即可；只要乘积不是指数级衰减，根式极限通常为 1。}
\end{solutionblock}

\begin{problemblock}
\textbf{例12.} 设
\[
a_n=\int_0^1x^n\sqrt{1-x^2}\,dx\qquad(n=0,1,2,\dots).
\]
(I) 证明：数列 \(\{a_n\}\) 单调递减，且
\[
a_n=\frac{n-1}{n+2}a_{n-2}\qquad(n=2,3,\dots);
\]
(II) 求
\[
\lim_{n\to\infty}\frac{a_n}{a_{n-1}}.
\]
\end{problemblock}

\begin{solutionblock}
\analysis{积分型数列先看被积函数随 \(n\) 的变化。递推式通常来自分部积分。比值极限可用递推式和夹逼。}

因为 \(0\le x\le1\)，所以
\[
x^{n+1}\sqrt{1-x^2}\le x^n\sqrt{1-x^2},
\]
且不恒等，因此
\[
a_{n+1}<a_n.
\]
故 \(\{a_n\}\) 单调递减。

对递推式，令
\[
a_n=\int_0^1x^{n-1}\cdot x\sqrt{1-x^2}\,dx.
\]
分部积分，取
\[
u=x^{n-1},\qquad dv=x\sqrt{1-x^2}\,dx.
\]
则
\[
du=(n-1)x^{n-2}\,dx,\qquad
v=-\frac13(1-x^2)^{3/2}.
\]
边界项为 0，故
\[
a_n=\frac{n-1}{3}\int_0^1x^{n-2}(1-x^2)^{3/2}\,dx.
\]
而
\[
(1-x^2)^{3/2}=(1-x^2)\sqrt{1-x^2},
\]
所以
\[
a_n=\frac{n-1}{3}\left(a_{n-2}-a_n\right).
\]
整理得
\[
(n+2)a_n=(n-1)a_{n-2},
\]
即
\[
a_n=\frac{n-1}{n+2}a_{n-2}.
\]

设
\[
r_n=\frac{a_n}{a_{n-1}}.
\]
由递推式
\[
a_n=\frac{n-1}{n+2}a_{n-2},\qquad
a_{n-1}=\frac{n-2}{n+1}a_{n-3},
\]
可得
\[
r_n=\frac{n-1}{n+2}\cdot\frac{a_{n-2}}{a_{n-1}}.
\]
由于 \(a_n\) 单调递减，\(0<r_n<1\)。又递推式说明相隔两项的比值趋于 1：
\[
\frac{a_n}{a_{n-2}}=\frac{n-1}{n+2}\to1.
\]
因此
\[
\frac{a_n}{a_{n-1}}\cdot\frac{a_{n-1}}{a_{n-2}}\to1.
\]
两个因子都在 \((0,1)\) 内，若其中一个远离 1，乘积不可能趋于 1，故
\[
\lim_{n\to\infty}\frac{a_n}{a_{n-1}}=1.
\]

\answer{\(\displaystyle \lim_{n\to\infty}\frac{a_n}{a_{n-1}}=1\)。}
\examnote{积分型数列中，\(a_n/a_{n-2}\to1\) 往往能推出相邻比值也趋于 1，但要补一句单调性保证比值夹在 \(0\) 与 \(1\) 之间。}
\end{solutionblock}

\begin{problemblock}
\textbf{例13.} 求极限
\[
\lim_{n\to\infty}\left[
\frac{\sin\frac{\pi}{n}}{n+1}
+\frac{\sin\frac{2\pi}{n}}{n+\frac12}
+\cdots+
\frac{\sin\pi}{n+\frac1n}
\right].
\]
\end{problemblock}

\begin{solutionblock}
\analysis{这是 Riemann 和型极限。分母 \(n+\frac1k\) 与 \(n\) 等价，核心和式近似为 \(\frac1n\sum_{k=1}^n\sin\frac{k\pi}{n}\)。}

一般项可写为
\[
\frac{\sin\frac{k\pi}{n}}{n+\frac1k}
=\frac1n\cdot
\frac{\sin\frac{k\pi}{n}}{1+\frac1{kn}}.
\]
因为
\[
0<\frac1{kn}\le1,
\]
且对除 \(k=1\) 外的大多数项有 \(\frac1{kn}\to0\)，整体可按 Riemann 和处理。更严格地，
\[
0\le
\frac1n\sin\frac{k\pi}{n}
-\frac{\sin\frac{k\pi}{n}}{n+\frac1k}
=\frac1n\sin\frac{k\pi}{n}\cdot
\frac{\frac1{kn}}{1+\frac1{kn}}
\le \frac1{n^2k}.
\]
求和得误差不超过
\[
\sum_{k=1}^n\frac1{n^2k}\le\frac{\ln n+1}{n^2}\to0.
\]
因此原极限等于
\[
\lim_{n\to\infty}\frac1n\sum_{k=1}^n\sin\frac{k\pi}{n}
=\int_0^1\sin(\pi x)\,dx
=\frac2\pi.
\]

\answer{\(\displaystyle \frac2\pi\)。}
\examnote{Riemann 和题要先把每项整理成 \(\frac1n f(k/n)\)，多出来的小扰动要证明总误差趋零。}
\end{solutionblock}

\begin{problemblock}
\textbf{例14.} (1) 设 \(k\) 为正常数，求
\[
\lim_{n\to\infty}
\frac{\ln n}{\ln(1+2^k+3^k+\cdots+n^k)}.
\]
(2) 求
\[
\lim_{n\to\infty}
\frac{(1^{2025}+2^{2025}+\cdots+n^{2025})^{2025}}
{(1^{2024}+2^{2024}+\cdots+n^{2024})^{2026}}.
\]
\end{problemblock}

\begin{solutionblock}
\analysis{本题使用幂和渐近公式：
\[
\sum_{j=1}^n j^p\sim\frac{n^{p+1}}{p+1}\qquad(p>-1).
\]
}

(1) 因为
\[
1+2^k+\cdots+n^k\sim\frac{n^{k+1}}{k+1},
\]
所以
\[
\ln(1+2^k+\cdots+n^k)
=\ln\left(\frac{n^{k+1}}{k+1}(1+o(1))\right)
=(k+1)\ln n+O(1).
\]
因此
\[
\lim_{n\to\infty}
\frac{\ln n}{\ln(1+2^k+\cdots+n^k)}
=\frac1{k+1}.
\]

(2) 由
\[
\sum_{j=1}^n j^{2025}\sim\frac{n^{2026}}{2026},\qquad
\sum_{j=1}^n j^{2024}\sim\frac{n^{2025}}{2025},
\]
得
\[
\frac{(1^{2025}+\cdots+n^{2025})^{2025}}
{(1^{2024}+\cdots+n^{2024})^{2026}}
\sim
\frac{\left(\frac{n^{2026}}{2026}\right)^{2025}}
{\left(\frac{n^{2025}}{2025}\right)^{2026}}.
\]
指数部分
\[
n^{2026\cdot2025-2025\cdot2026}=1,
\]
故极限为
\[
\frac{2025^{2026}}{2026^{2025}}.
\]

\answer{(1) \(\displaystyle \frac1{k+1}\)；(2) \(\displaystyle \frac{2025^{2026}}{2026^{2025}}\)。}
\examnote{幂和只需记住最高阶：\(\sum_{j=1}^n j^p\) 的阶是 \(n^{p+1}\)，常数是 \(\frac1{p+1}\)。}
\end{solutionblock}

\begin{problemblock}
\textbf{例15.} 数列 \(\{a_n\},\{b_n\}\) 满足
\[
b_{n-1}\le a_n\le b_n.
\]
“\(\{a_n\}\) 收敛”是“\(\{b_n\}\) 收敛”的（\quad）

A. 充分不必要条件 \qquad B. 必要不充分条件

C. 既不必要又不充分条件 \qquad D. 充要条件
\end{problemblock}

\begin{solutionblock}
\analysis{关键是把不等式错位使用。由 \(b_{n-1}\le a_n\le b_n\)，把 \(n\) 换成 \(n+1\)，还可得 \(b_n\le a_{n+1}\)。}

若 \(\{b_n\}\) 收敛，设 \(b_n\to L\)，则
\[
b_{n-1}\to L,\qquad b_n\to L.
\]
由
\[
b_{n-1}\le a_n\le b_n
\]
夹逼得
\[
a_n\to L.
\]
所以 \(\{b_n\}\) 收敛推出 \(\{a_n\}\) 收敛。

反过来，若 \(\{a_n\}\) 收敛，设 \(a_n\to A\)。由原式和错位式
\[
a_n\le b_n\le a_{n+1},
\]
夹逼得
\[
b_n\to A.
\]
所以 \(\{a_n\}\) 收敛也推出 \(\{b_n\}\) 收敛。

因此二者互相推出。

\answer{D}
\examnote{遇到 \(b_{n-1},a_n,b_n\) 这种错位不等式，要同时写出 \(a_n\le b_n\le a_{n+1}\)。}
\end{solutionblock}

\begin{problemblock}
\textbf{例16.} 设
\[
x_1=\ln a\quad(a>0),\qquad
x_{n+1}=x_n+\ln(a-x_n).
\]
证明数列 \(\{x_n\}\) 收敛，并求
\[
\lim_{n\to\infty}x_n.
\]
\end{problemblock}

\begin{solutionblock}
\analysis{固定点由 \(L=L+\ln(a-L)\) 给出，即 \(L=a-1\)。证明时把变量平移到这个固定点。}

令
\[
y_n=x_n-(a-1).
\]
则
\[
a-x_n=1-y_n,
\]
且
\[
y_{n+1}=y_n+\ln(1-y_n).
\]
初值
\[
y_1=\ln a-a+1\le0
\]
由基本不等式 \(\ln a\le a-1\) 得到。

若 \(y_n\le0\)，则 \(1-y_n\ge1\)，所以
\[
\ln(1-y_n)\ge0,
\]
从而
\[
y_{n+1}\ge y_n.
\]
又函数
\[
\psi(y)=y+\ln(1-y)
\]
在 \((-\infty,0]\) 上满足 \(\psi(y)\le0\)，因为 \(\psi(0)=0\)，且
\[
\psi'(y)=1-\frac1{1-y}=\frac{-y}{1-y}\ge0.
\]
于是 \(y_n\le0\) 对所有 \(n\) 成立，且 \(\{y_n\}\) 单调递增有上界 0，故收敛。

设 \(y_n\to l\le0\)，由递推式取极限：
\[
l=l+\ln(1-l),
\]
故 \(l=0\)。因此
\[
\lim_{n\to\infty}x_n=a-1.
\]

\answer{\(\displaystyle a-1\)。}
\examnote{抽象递推先找不动点，再作平移，通常能把问题化成单调有界。}
\end{solutionblock}

\begin{problemblock}
\textbf{例17.} 设
\[
x_1=2,\qquad x_{n+1}=\frac{x_n^2}{1+x_n}\quad(n=1,2,\dots).
\]
(1) 证明数列 \(\{x_n\}\) 收敛，并求 \(\displaystyle\lim_{n\to\infty}x_n\)；

(2) 求
\[
\lim_{n\to\infty}\left(
\frac{x_2}{x_1}+\frac{x_3}{x_2}+\cdots+\frac{x_{n+1}}{x_n}
\right).
\]
\end{problemblock}

\begin{solutionblock}
\analysis{第一问单调有界。第二问的关键是发现比值可以化成相邻项差。}

因为 \(x_1=2>0\)，若 \(x_n>0\)，则 \(x_{n+1}>0\)，且
\[
x_{n+1}=\frac{x_n^2}{1+x_n}<x_n.
\]
所以 \(\{x_n\}\) 单调递减且有下界 0，故收敛。设极限为 \(L\ge0\)，则
\[
L=\frac{L^2}{1+L}.
\]
即
\[
L(1+L)=L^2,
\]
故 \(L=0\)。

第二问中
\[
\frac{x_{k+1}}{x_k}=\frac{x_k}{1+x_k}.
\]
另一方面
\[
x_k-x_{k+1}
=x_k-\frac{x_k^2}{1+x_k}
=\frac{x_k}{1+x_k}.
\]
所以
\[
\frac{x_{k+1}}{x_k}=x_k-x_{k+1}.
\]
于是
\[
\sum_{k=1}^n\frac{x_{k+1}}{x_k}
=\sum_{k=1}^n(x_k-x_{k+1})
=x_1-x_{n+1}.
\]
令 \(n\to\infty\)，得
\[
\lim_{n\to\infty}\sum_{k=1}^n\frac{x_{k+1}}{x_k}
=2-0=2.
\]

\answer{(1) 极限为 \(0\)；(2) 极限为 \(2\)。}
\examnote{递推式里出现 \(\frac{x_n^2}{1+x_n}\)，要试着计算 \(x_n-x_{n+1}\)，常能得到望远镜结构。}
\end{solutionblock}

\begin{problemblock}
\textbf{例18.} 设 \(0<x_1<1\)，
\[
x_n=\int_0^1\max(x_{n-1},t)\,dt\qquad(n=2,3,\dots).
\]
证明数列 \(\{x_n\}\) 收敛并求极限。
\end{problemblock}

\begin{solutionblock}
\analysis{先把积分递推化成显式递推。由于 \(0<x_{n-1}<1\)，积分区间可按 \(t=x_{n-1}\) 分段。}

若 \(0<c<1\)，则
\[
\int_0^1\max(c,t)\,dt
=\int_0^c c\,dt+\int_c^1t\,dt
=c^2+\frac{1-c^2}{2}
=\frac{1+c^2}{2}.
\]
因此
\[
x_n=\frac{1+x_{n-1}^2}{2}.
\]
若 \(0<x_{n-1}<1\)，则
\[
0<x_n<1,
\]
且
\[
x_n-x_{n-1}
=\frac{1+x_{n-1}^2-2x_{n-1}}2
=\frac{(1-x_{n-1})^2}{2}>0.
\]
所以 \(\{x_n\}\) 单调递增且有上界 1，故收敛。设极限为 \(L\)，则
\[
L=\frac{1+L^2}{2},
\]
即
\[
(L-1)^2=0.
\]
因此
\[
L=1.
\]

\answer{\(\displaystyle 1\)。}
\examnote{含 \(\max\) 的积分递推一般先按最大值切换点分段，把递推式显式化。}
\end{solutionblock}

\begin{problemblock}
\textbf{例19.} 设正项数列 \(\{x_n\}\) 满足
\[
x_1=\sqrt2,\qquad x_{n+1}^2=2^{x_n}\quad(n=1,2,\dots).
\]
证明 \(\{x_n\}\) 收敛，并求
\[
\lim_{n\to\infty}x_n.
\]
\end{problemblock}

\begin{solutionblock}
\analysis{递推可写成 \(x_{n+1}=2^{x_n/2}\)。固定点有 \(2\) 和 \(4\)，但初值在 \((0,2)\)，因此极限只能落在 2。}

先证明
\[
\sqrt2\le x_n<2.
\]
若 \(0<x_n<2\)，则
\[
x_{n+1}=2^{x_n/2}<2.
\]
又函数 \(2^{x/2}\) 在 \([\sqrt2,2]\) 上大于等于 \(\sqrt2\)，故区间不变。

再证明单调增加。令
\[
g(x)=2^{x/2}-x.
\]
在 \((0,2)\) 内可验证 \(g(x)>0\)，例如 \(g(\sqrt2)>0\)，且 \(x=2\) 是该区间右端固定点。因此
\[
x_{n+1}>x_n.
\]
于是 \(\{x_n\}\) 单调递增且有上界 2，故收敛。

设极限为 \(L\in[\sqrt2,2]\)，由递推得
\[
L^2=2^L.
\]
在 \([\sqrt2,2]\) 内唯一解为
\[
L=2.
\]

\answer{\(\displaystyle 2\)。}
\examnote{递推 \(x_{n+1}=\varphi(x_n)\) 的题，先找不变区间，再证单调，最后解固定点。}
\end{solutionblock}

\begin{problemblock}
\textbf{例20.} 设 \(f(x)\) 是区间 \([0,+\infty)\) 上单调减少且非负的连续函数，
\[
a_n=\sum_{k=1}^n f(k)-\int_1^n f(x)\,dx\qquad(n=1,2,\dots).
\]
证明：数列 \(\{a_n\}\) 的极限存在。
\end{problemblock}

\begin{solutionblock}
\analysis{要证极限存在，常用单调有界。这里利用单调减少函数的积分夹估。}

先看差分：
\[
a_{n+1}-a_n
=f(n+1)-\int_n^{n+1}f(x)\,dx.
\]
由于 \(f\) 单调减少，所以当 \(x\in[n,n+1]\) 时，
\[
f(n+1)\le f(x)\le f(n).
\]
因此
\[
\int_n^{n+1}f(x)\,dx\ge f(n+1),
\]
从而
\[
a_{n+1}-a_n\le0.
\]
故 \(\{a_n\}\) 单调不增。

再证有下界。由同样的单调性，
\[
\int_k^{k+1}f(x)\,dx\le f(k).
\]
所以
\[
\int_1^n f(x)\,dx
=\sum_{k=1}^{n-1}\int_k^{k+1}f(x)\,dx
\le \sum_{k=1}^{n-1}f(k).
\]
于是
\[
a_n
=\sum_{k=1}^n f(k)-\int_1^n f(x)\,dx
\ge f(n)\ge0.
\]
因此 \(\{a_n\}\) 单调不增且有下界 0，必收敛。

\answer{数列 \(\{a_n\}\) 单调不增且有下界，因此 \(\displaystyle\lim_{n\to\infty}a_n\) 存在。}
\method{方法二（Cauchy 判别法）}
也可直接证 Cauchy。对 \(m>n\)，利用单调性把 \(a_n-a_m\) 写成若干小区间上“左端矩形面积减积分”的和，并由 \(0\le a_n-a_m\le f(n)\) 型估计控制尾差；因为 \(f(n)\) 有极限且非负，此法同样能推出 \(a_n\) 收敛。
\examnote{“积分与和式比较”是考研常考模型：单调函数在每个小区间上的最大值、最小值给出上下估计。}
\end{solutionblock}

\begin{problemblock}
\textbf{例21.} 设
\[
I_n=\int_0^{\pi/2}\frac{\sin^2 nt}{\sin t}\,dt,
\]
其中 \(n\) 是正整数，证明：
\[
(1)\ I_n-I_{n-1}=\frac1{2n-1};
\]
\[
(2)\ \lim_{n\to\infty}(2I_n-\ln n)\text{ 存在}.
\]
\end{problemblock}

\begin{solutionblock}
\analysis{第一问用三角恒等式化差；第二问把 \(I_n\) 写成奇数调和和，再与调和级数比较。}

由
\[
\sin^2A-\sin^2B=\sin(A+B)\sin(A-B),
\]
取 \(A=nt,\ B=(n-1)t\)，得
\[
\sin^2nt-\sin^2(n-1)t=\sin(2n-1)t\sin t.
\]
因此
\[
I_n-I_{n-1}
=\int_0^{\pi/2}\sin(2n-1)t\,dt
=\frac{1-\cos\frac{(2n-1)\pi}{2}}{2n-1}
=\frac1{2n-1}.
\]

又 \(I_0=0\)，所以
\[
I_n=\sum_{k=1}^n\frac1{2k-1}.
\]
于是
\[
2I_n=2\sum_{k=1}^n\frac1{2k-1}
=2H_{2n}-H_n,
\]
其中 \(H_n=\sum_{k=1}^n\frac1k\)。由
\[
H_n=\ln n+\gamma+o(1)
\]
可得
\[
2I_n-\ln n
=2H_{2n}-H_n-\ln n
\to 2\ln2+\gamma.
\]
所以极限存在。

\answer{(1) \(I_n-I_{n-1}=\dfrac1{2n-1}\)；(2) \(\displaystyle\lim_{n\to\infty}(2I_n-\ln n)=2\ln2+\gamma\)，故极限存在。}
\method{方法二（单调有界法）}
多解补充：也可以不使用调和数渐近，而证明 \(2I_n-\ln n\) 单调有界。由 \(I_n=\sum_{k=1}^n\frac1{2k-1}\)，把 \(\ln n\) 写成 \(\sum_{k=1}^{n-1}\ln(1+1/k)\)，再用 \(\frac1{k+1}<\ln(1+1/k)<\frac1k\) 夹估，即可得到收敛性；若要求精确极限，再回到 \(H_n=\ln n+\gamma+o(1)\)。
\examnote{含 \(\sin^2nt-\sin^2(n-1)t\) 的差分，优先用平方差三角公式；它会直接消掉分母 \(\sin t\)。}
\end{solutionblock}

\begin{problemblock}
\textbf{例22.} 设正项数列 \(\{x_n\}\) 满足
\[
x_{n+1}\le x_n+\frac{2^n}{(2^{n+1}-1)^2}\qquad(n=1,2,\dots),
\]
则下列说法正确的有几项？
\[
\begin{array}{ll}
\text{① }\left\{x_n+\dfrac1{2^n-1}\right\}\text{ 递减；}&
\text{② }\{x_n\}\text{ 收敛；}\\[2mm]
\text{③ }\displaystyle\sum_{k=1}^n x_k<(x_1+1)(n+1);&
\text{④ }\exists N_1\in\mathbb N^+,\ n>N_1\text{ 时 }x_n<1+\dfrac{x_1}{2}.
\end{array}
\]
A. 1 \qquad B. 2 \qquad C. 3 \qquad D. 4
\end{problemblock}

\begin{solutionblock}
\analysis{这题的关键是看出给定增量小于一个可望远镜的差：
\[
\frac1{2^n-1}-\frac1{2^{n+1}-1}.
\]
}

因为
\[
\frac1{2^n-1}-\frac1{2^{n+1}-1}
=\frac{2^n}{(2^n-1)(2^{n+1}-1)}
>
\frac{2^n}{(2^{n+1}-1)^2},
\]
所以
\[
x_{n+1}-x_n
\le \frac{2^n}{(2^{n+1}-1)^2}
<
\frac1{2^n-1}-\frac1{2^{n+1}-1}.
\]
移项得
\[
x_{n+1}+\frac1{2^{n+1}-1}
<
x_n+\frac1{2^n-1}.
\]
故①正确。

由①知
\[
y_n=x_n+\frac1{2^n-1}
\]
单调递减。又 \(x_n>0\)，所以 \(y_n>0\)，因此 \(y_n\) 收敛。由于
\[
\frac1{2^n-1}\to0,
\]
可得 \(x_n\) 收敛，②正确。

由①还得
\[
x_n<x_n+\frac1{2^n-1}\le x_1+1,
\]
因此
\[
\sum_{k=1}^n x_k<n(x_1+1)<(n+1)(x_1+1),
\]
③正确。

④不一定正确。例如取常数列 \(x_n\equiv x_1\)，它满足题设不等式。若 \(x_1\ge2\)，则
\[
x_n=x_1\not<1+\frac{x_1}{2}.
\]
故④错误。

正确的是①②③，共 3 项。

\answer{C}
\examnote{数列不等式题常把右侧增量改写成“某个辅助项的差”，从而构造单调量。}
\end{solutionblock}

\begin{problemblock}
\textbf{例23.} (1) 证明方程
\[
\ln x+x^n=0
\]
在 \((0,1)\) 内有唯一实根 \(x_n\) \((n=1,2,\dots)\)；

(2) 证明 \(\displaystyle\lim_{n\to\infty}x_n\) 存在。
\end{problemblock}

\begin{solutionblock}
\analysis{先用单调性证明根唯一，再用任意 \(c<1\) 的定位证明根趋于 1。}

令
\[
F_n(x)=\ln x+x^n,\qquad x\in(0,1].
\]
则
\[
F_n'(x)=\frac1x+nx^{n-1}>0,
\]
所以 \(F_n\) 严格递增。又
\[
\lim_{x\to0^+}F_n(x)=-\infty,\qquad F_n(1)=1>0,
\]
故在 \((0,1)\) 内有唯一实根 \(x_n\)。

下面证明 \(x_n\to1\)。任取 \(c\in(0,1)\)，则
\[
F_n(c)=\ln c+c^n\to\ln c<0.
\]
所以充分大时 \(F_n(c)<0\)。又 \(F_n(x_n)=0\)，且 \(F_n\) 递增，因此充分大时
\[
x_n>c.
\]
由于 \(x_n<1\)，任意 \(c<1\) 最终都有 \(c<x_n<1\)，故
\[
\lim_{n\to\infty}x_n=1.
\]

\answer{极限存在且为 \(1\)。}
\examnote{隐式根数列常用“固定 \(c\)，比较 \(F_n(c)\) 的符号”来定位极限。}
\end{solutionblock}

\begin{problemblock}
\textbf{例24.} 设
\[
f_n(x)=\sin x+\sin^2x+\cdots+\sin^n x\qquad(n=1,2,\dots).
\]
(1) 证明方程 \(f_n(x)=1\) 在 \(\left[\frac\pi6,\frac\pi2\right]\) 上有且只有一个实根 \(x_n\)；

(2) 求
\[
\lim_{n\to\infty}x_n.
\]
\end{problemblock}

\begin{solutionblock}
\analysis{把 \(\sin x\) 看成新变量 \(y\)。方程是有限等比和 \(y+y^2+\cdots+y^n=1\)。}

在 \(\left[\frac\pi6,\frac\pi2\right]\) 上，\(\sin x\) 单调增加，因此 \(f_n(x)\) 严格增加。并且
\[
f_n\left(\frac\pi6\right)
=\frac12+\frac1{2^2}+\cdots+\frac1{2^n}
=1-\frac1{2^n}<1,
\]
而
\[
f_n\left(\frac\pi2\right)=n\ge1,
\]
当 \(n\ge2\) 时大于 1；\(n=1\) 时根为 \(\pi/2\)。故存在唯一根。

设
\[
y_n=\sin x_n.
\]
则
\[
y_n+y_n^2+\cdots+y_n^n=1,\qquad \frac12\le y_n\le1.
\]
若 \(y_n\) 有子列趋于 \(y\in[\frac12,1]\)。若 \(y>\frac12\)，则
\[
y+y^2+\cdots
=\frac{y}{1-y}>1,
\]
与有限和等于 1 的极限趋势矛盾；若 \(y=\frac12\)，则正好对应极限和 1。故
\[
y_n\to\frac12.
\]
于是
\[
x_n\to\arcsin\frac12=\frac\pi6.
\]

\answer{\(\displaystyle \frac\pi6\)。}
\examnote{有限和方程的根极限，通常先猜无限级数方程，再证明根不会跑到别处。}
\end{solutionblock}

\begin{problemblock}
\textbf{例25.} 设 \(f(x)\) 是 \([0,1]\) 上的正值连续函数，证明：

(1) 存在唯一的 \(\xi\in(0,1)\)，使得
\[
\int_0^\xi f(t)\,dt=\int_\xi^1\frac1{f(t)}\,dt.
\]

(2) 对任意自然数 \(n\ge2\)，存在唯一的 \(x_n\in(0,1)\)，使得
\[
\int_{1/n}^{x_n} f(t)\,dt=\int_{x_n}^1\frac1{f(t)}\,dt,
\]
且
\[
\lim_{n\to\infty}x_n=\xi.
\]
\end{problemblock}

\begin{solutionblock}
\analysis{把等式移到一边，构造严格单调连续函数。第二问是方程随参数 \(1/n\to0\) 的根连续性。}

定义
\[
F(x)=\int_0^x f(t)\,dt-\int_x^1\frac1{f(t)}\,dt.
\]
因为 \(f>0\) 且连续，
\[
F'(x)=f(x)+\frac1{f(x)}>0.
\]
所以 \(F\) 严格递增。又
\[
F(0)=-\int_0^1\frac1{f(t)}\,dt<0,\qquad
F(1)=\int_0^1f(t)\,dt>0.
\]
由介值定理和严格单调性，存在唯一 \(\xi\in(0,1)\)，使 \(F(\xi)=0\)。

对第二问，定义
\[
F_n(x)=\int_{1/n}^x f(t)\,dt-\int_x^1\frac1{f(t)}\,dt.
\]
同样有
\[
F_n'(x)=f(x)+\frac1{f(x)}>0.
\]
且
\[
F_n(1/n)<0,\qquad F_n(1)>0,
\]
故存在唯一 \(x_n\in(1/n,1)\subset(0,1)\)，使 \(F_n(x_n)=0\)。

注意
\[
F_n(x)=F(x)-\int_0^{1/n}f(t)\,dt.
\]
由 \(F_n(x_n)=0\)，得
\[
F(x_n)=\int_0^{1/n}f(t)\,dt\to0.
\]
由于 \(F\) 连续严格递增，且 \(F(\xi)=0\)，故
\[
x_n\to\xi.
\]

\answer{(1) 唯一的 \(\xi\in(0,1)\) 存在；(2) 对每个 \(n\ge2\)，唯一的 \(x_n\in(0,1)\) 存在，且 \(\displaystyle\lim_{n\to\infty}x_n=\xi\)。}
\method{方法二（反证定位法）}
也可不用一致收敛语言。任取 \(\varepsilon>0\)，由 \(F\) 严格递增知 \(F(\xi-\varepsilon)<0<F(\xi+\varepsilon)\)。当 \(n\) 充分大时，\(\int_0^{1/n}f(t)\,dt\) 足够小，于是 \(F_n(\xi-\varepsilon)<0<F_n(\xi+\varepsilon)\)，唯一零点 \(x_n\) 被夹在 \(\xi-\varepsilon\) 与 \(\xi+\varepsilon\) 之间，从而 \(x_n\to\xi\)。
\examnote{存在唯一性题的标准模板：连续保证存在，严格单调保证唯一。根的极限用方程函数一致逼近或直接代回。}
\end{solutionblock}

\begin{problemblock}
\textbf{例26.} (1) 证明方程
\[
e^x+x^{2n+1}=0
\]
在 \((-1,0)\) 内有唯一实根 \(x_n\) \((n=0,1,2,\dots)\)；

(2) 证明 \(\displaystyle\lim_{n\to\infty}x_n\) 存在并求其值 \(a\)；

(3) 求
\[
\lim_{n\to\infty}n(x_n-a).
\]
\end{problemblock}

\begin{solutionblock}
\analysis{根在 \((-1,0)\)。当 \(n\to\infty\) 时，若 \(x\) 固定在 \((-1,0)\)，则 \(x^{2n+1}\to0\)，所以根只能逼近 \(-1\)。第三问要对 \(x_n=-1+u_n\) 取对数。}

令
\[
F_n(x)=e^x+x^{2n+1}.
\]
则
\[
F_n' (x)=e^x+(2n+1)x^{2n}>0.
\]
又
\[
F_n(-1)=e^{-1}-1<0,\qquad F_n(0)=1>0.
\]
故在 \((-1,0)\) 内有唯一实根 \(x_n\)。

对任意 \(c\in(-1,0)\)，有
\[
F_n(c)=e^c+c^{2n+1}\to e^c>0.
\]
而 \(F_n(-1)<0\)，根 \(x_n\) 必最终落在 \((-1,c)\) 内。任意 \(c>-1\) 都成立，故
\[
x_n\to -1.
\]
于是
\[
a=-1.
\]

令
\[
x_n=-1+u_n,\qquad u_n\to0^+.
\]
由方程
\[
e^{-1+u_n}+(-1+u_n)^{2n+1}=0
\]
得
\[
(1-u_n)^{2n+1}=e^{-1+u_n}.
\]
取对数：
\[
(2n+1)\ln(1-u_n)=-1+u_n.
\]
因为
\[
\ln(1-u_n)=-u_n+O(u_n^2),
\]
故
\[
-(2n+1)u_n+o(nu_n)=-1+u_n.
\]
由此先知 \(nu_n\) 有界，再令极限得
\[
2nu_n\to1.
\]
所以
\[
\lim_{n\to\infty}n(x_n-a)
=\lim_{n\to\infty}n(x_n+1)
=\frac12.
\]

\answer{\(a=-1\)，且 \(\displaystyle\lim_{n\to\infty}n(x_n-a)=\frac12\)。}
\examnote{根靠近端点时，令 \(x_n=-1+u_n\) 并取对数，是处理指数幂次平衡的常用方法。}
\end{solutionblock}

\begin{problemblock}
\textbf{例27.} (1) 数列 \(\{x_n\}\) 满足
\[
x_{n+1}=f(x_n),
\]
\(f(x)\) 一阶可导。若 \(a\) 是 \(f(x)=x\) 的唯一解，且存在常数 \(k\in(0,1)\)，使
\[
|f'(x)|<k<1,
\]
证明 \(\{x_n\}\) 收敛于 \(a\)。

(2) 设
\[
x_1=\frac a2\quad(0<a\le1),\qquad
x_{n+1}=\frac a2-\frac{x_n^2}{2}.
\]
求 \(\displaystyle\lim_{n\to\infty}x_n\)。
\end{problemblock}

\begin{solutionblock}
\analysis{第一问是压缩映射思想。第二问可直接把 \(f(x)=\frac a2-\frac{x^2}{2}\) 放入第一问框架。}

(1) 因为 \(a=f(a)\)，由中值定理，
\[
|x_{n+1}-a|
=|f(x_n)-f(a)|
\le k|x_n-a|.
\]
递推得
\[
|x_n-a|\le k^{n-1}|x_1-a|\to0.
\]
所以
\[
x_n\to a.
\]

(2) 设
\[
f(x)=\frac a2-\frac{x^2}{2}.
\]
在区间 \([0,a/2]\) 上，
\[
|f'(x)|=|x|\le\frac a2\le\frac12<1,
\]
且 \(f\) 把该区间映入自身。因此数列收敛于 \(f(x)=x\) 的唯一解。

解
\[
x=\frac a2-\frac{x^2}{2}
\]
得
\[
x^2+2x-a=0,
\]
所以
\[
x=-1\pm\sqrt{1+a}.
\]
由于数列非负，极限为
\[
\sqrt{1+a}-1.
\]

\answer{\(\displaystyle \sqrt{1+a}-1\)。}
\examnote{若递推函数导数绝对值小于 1，就优先考虑压缩估计；这比硬证单调更省力。}
\end{solutionblock}

\begin{problemblock}
\textbf{例28.} 设
\[
-1\le x_1\le1,\qquad x_{n+1}=\cos x_n.
\]
证明
\[
\lim_{n\to\infty}x_n
\]
存在。
\end{problemblock}

\begin{solutionblock}
\analysis{这是经典不动点迭代。虽然 \(\cos x\) 不一定使数列单调，但在 \([-1,1]\) 上它是压缩映射。}

因为 \(x_1\in[-1,1]\)，且 \(\cos([-1,1])\subset[\cos1,1]\subset[-1,1]\)，所以所有 \(x_n\in[-1,1]\)。

对任意 \(u,v\in[-1,1]\)，由中值定理
\[
|\cos u-\cos v|\le(\sin1)|u-v|.
\]
注意 \(\sin1<1\)。设 \(a\) 是方程
\[
a=\cos a
\]
在 \([-1,1]\) 内的唯一解，则
\[
|x_{n+1}-a|
=|\cos x_n-\cos a|
\le(\sin1)|x_n-a|.
\]
于是
\[
|x_n-a|\le(\sin1)^{n-1}|x_1-a|\to0.
\]
故 \(x_n\to a\)，极限存在。

\answer{\(\displaystyle\lim_{n\to\infty}x_n=a\)，其中 \(a\) 是方程 \(a=\cos a\) 在 \([-1,1]\) 内的唯一实根。}
\method{方法二（偶奇子列法）}
还可以把迭代两次：\(x_{n+2}=\cos(\cos x_n)\)。函数 \(g(x)=\cos(\cos x)\) 在 \([0,1]\) 上递增，且 \(g(x)-x\) 只有一个零点。分别研究偶子列和奇子列，可得二者收敛到同一不动点 \(a\)，于是原数列收敛。考场上压缩映射更短。
\examnote{\(x_{n+1}=\cos x_n\) 不适合直接证单调；用压缩映射最干净。}
\end{solutionblock}

\begin{problemblock}
\textbf{例29.} 设
\[
x_1=2,\qquad x_{n+1}=2+\frac1{x_n}.
\]
求证 \(\displaystyle\lim_{n\to\infty}x_n\) 存在并求其值。
\end{problemblock}

\begin{solutionblock}
\analysis{递推函数 \(f(x)=2+\frac1x\) 在 \([2,3]\) 上是压缩映射。}

先由 \(x_1=2\) 得 \(x_2=5/2\)，且若 \(2\le x_n\le3\)，则
\[
2<x_{n+1}=2+\frac1{x_n}\le \frac52<3.
\]
所以 \(x_n\in[2,3]\)。

在 \([2,3]\) 上，
\[
|f'(x)|=\frac1{x^2}\le\frac14<1.
\]
因此迭代收敛到唯一不动点 \(L\)。解
\[
L=2+\frac1L
\]
得
\[
L^2-2L-1=0.
\]
由于 \(L>0\)，
\[
L=1+\sqrt2.
\]

\answer{\(\displaystyle 1+\sqrt2\)。}
\examnote{分式递推常先找不变区间，再用导数小于 1 证明收敛。}
\end{solutionblock}

\begin{problemblock}
\textbf{例30.} 设
\[
x_1>x_2>0,\qquad x_{n+2}=\sqrt{x_{n+1}x_n}.
\]
证明
\[
\lim_{n\to\infty}x_n
\]
存在，并求其值。
\end{problemblock}

\begin{solutionblock}
\analysis{几何平均递推取对数后变为线性递推，这是本题关键。}

令
\[
y_n=\ln x_n.
\]
则
\[
y_{n+2}=\frac{y_{n+1}+y_n}{2}.
\]
设
\[
d_n=y_{n+1}-y_n.
\]
则
\[
d_{n+1}=y_{n+2}-y_{n+1}
=-\frac12(y_{n+1}-y_n)
=-\frac12d_n.
\]
因此
\[
d_n=\left(-\frac12\right)^{n-1}d_1.
\]
于是
\[
y_n=y_1+\sum_{k=1}^{n-1}d_k
\to y_1+\sum_{k=1}^{\infty}\left(-\frac12\right)^{k-1}(y_2-y_1).
\]
几何级数求和得
\[
\lim_{n\to\infty}y_n
=y_1+\frac{1}{1+\frac12}(y_2-y_1)
=\frac13y_1+\frac23y_2.
\]
所以
\[
\lim_{n\to\infty}x_n
=\exp\left(\frac13\ln x_1+\frac23\ln x_2\right)
=\sqrt[3]{x_1x_2^2}.
\]

\answer{\(\displaystyle \sqrt[3]{x_1x_2^2}\)。}
\examnote{乘法、幂、几何平均递推，优先取对数转成加法线性递推。}
\end{solutionblock}

\begin{problemblock}
\textbf{例31.} 设数列 \(\{x_n\}\) 定义如下：
\[
x_1=1,\qquad x_n=x_{n+1}+e^{x_{n+1}}-1\quad(n=1,2,\dots).
\]
证明数列 \(\{x_n\}\) 收敛并求
\[
\lim_{n\to\infty}x_n.
\]
\end{problemblock}

\begin{solutionblock}
\analysis{与例6类似，这是反向定义递推。函数 \(t+e^t-1\) 严格递增，因此可唯一确定下一项。}

设
\[
\phi(t)=t+e^t-1.
\]
则
\[
\phi'(t)=1+e^t>0,
\]
所以 \(\phi\) 严格递增，且 \(\phi(0)=0\)。

由 \(x_1=1>0\)，若 \(x_n>0\)，方程
\[
\phi(x_{n+1})=x_n
\]
有唯一正解 \(x_{n+1}>0\)。又当 \(x_{n+1}>0\) 时，
\[
x_n=x_{n+1}+e^{x_{n+1}}-1>x_{n+1}.
\]
因此 \(\{x_n\}\) 单调递减且有下界 0，故收敛。

设极限为 \(L\ge0\)。对递推式取极限：
\[
L=L+e^L-1.
\]
故
\[
e^L=1,\qquad L=0.
\]

\answer{\(\displaystyle 0\)。}
\examnote{反向递推 \(x_n=\phi(x_{n+1})\) 常用 \(\phi\) 的单调性证明下一项存在唯一，再判断 \(x_{n+1}<x_n\)。}
\end{solutionblock}
